\section{Discussion}
\label{sec:discussion}

\paragraph{CIVO mechanism.}
\looseness=-1 CIVO propagates through multiple, environment-dependent pathways: Walker Walk develops a $1.63\times$ reward prediction gap while Cheetah Run shows none ($1.00\times$) despite higher severity ($8.27\times$ NVE). In Cheetah, observation fidelity reverses ($0.93\times$; $\dt{=}2$ predictions more accurate) while embedding geometry separates comparably (cosine ${\leq}{-}0.84$), suggesting CIVO propagates through latent-space dynamics rather than reward prediction. A random-policy control experiment ($N{=}100$ rollouts, 95\% bootstrap CI) decomposes the effect into two necessary stages: under a random policy, $\dt{=}2$ imagination produces $1.68\times$ higher mean reward than $\dt{=}1$ in the trained-embedding model, attributing ${\approx}36\%$ of overestimation to dynamics bias (the RSSM prior already generates reward-favorable trajectories at $\dt{>}1$); the trained actor amplifies this to $2.90\times$, attributing ${\approx}64\%$ to actor exploitation (in the no-embedding control, neither stage is present: random-policy ratio $1.11\times$, trained-actor ratio $0.98\times$; Appendix~\ref{sec:civo_mechanism}). Signed reward error analysis reinforces this: $\dt{=}2$ rewards are \emph{under}predicted (mean SE${=}{-}0.599$, only 8\% of samples overpredicted), ruling out observation-level reward overprediction and localizing CIVO to imagination-level prior drift---the RSSM prior generates reward-favorable trajectories unconstrained by observations. The mechanism is analogous to delusional bias \citep{lu2018nondelusional} but model-induced; unlike MPC, imagination-based actor-critic bootstraps value targets entirely within the learned model without real-data correction (Appendix~\ref{sec:mpc_immune}). Frozen embeddings participate identically in the forward pass yet never trigger CIVO (0/16 controls vs.\ 5/7 trained, $p < 0.001$; Table~\ref{tab:causal_ablation}): the vulnerability lies in what the embedding \emph{learns}, not in conditioning architecture.

\paragraph{Catastrophic consequence.}
\looseness=-1 Even at the lower bound of onset probability (Clopper-Pearson 95\% CI: Walker $[15\%, 95\%]$, Cheetah $[19\%, 99\%]$), every CIVO-positive seed suffers irrecoverable performance collapse (eval $< 30\%$ of baseline), making diagnosis essential regardless of onset frequency. Of the 5/7 multi-joint seeds that trigger CIVO, 3 sustain permanent collapse, 1 self-corrects with permanent performance deficit, and 1 shows NVE self-correction without policy recovery. The binary nature of this outcome (healthy \emph{or} collapsed, with no graceful degradation) means that standard training practices cannot tolerate even low onset probability: a single CIVO event wastes an entire training run. This argues for proactive monitoring (NVE as early warning) rather than post-hoc analysis.

\paragraph{Practical implications.}
\label{sec:detector}
\looseness=-1 The systematic resilience documented in \S\ref{sec:remedies} yields three practitioner guidelines. First, monitor NVE continuously during training: a sustained NVE $> 2.0\times$ after 35K steps is diagnostic of CIVO with AUROC${=}1.000$ ($n{=}36$), enabling early termination before compute is wasted. Second, validate with frozen-embedding controls before deploying learned temporal embeddings: frozen embeddings never trigger CIVO (0/16 controls) while preserving the forward-pass architecture, isolating whether observed degradation is embedding-specific. Third, prefer MPC-style re-planning for variable-rate execution: MPC re-plans from real states at each step, so model bias at $\dt{>}1$ is corrected by real observations before accumulating into value estimates, the structural factor that distinguishes immune MPC from vulnerable imagination-based actor-critic.

\paragraph{Implications beyond DreamerV3.}
\looseness=-1 The structural vulnerability---shared imagined rollouts defeat diversity-based remedies---is not architecture-specific. Any imagination-based actor-critic that (1)~trains its critic entirely on model-generated rollouts and (2)~bootstraps value targets within the same learned model faces the same closed-loop constraint: remedies that assume diverse or independent training sources (Double Q pessimism, ensemble disagreement, noise injection) lose corrective power when both sides of the training signal originate from the same model. Any future architecture that combines temporal conditioning with pure-imagination value estimation would face the same constraint---for instance, extending TD-MPC2 \citep{hansen2024tdmpc2} to train critics on imagination rather than real trajectories, or adding $\dt$-conditioning to transformer world models such as IRIS \citep{micheli2023iris}. Whether the specific onset boundary (6~DoF locomotion) generalizes remains open; the structural constraint does.

\paragraph{Environment complexity boundary.}
\looseness=-1 Cross-environment analysis across six DMC tasks spanning 1--21 DoF reveals that CIVO occupies a complexity Goldilocks zone: 1-DoF (Cartpole, 0/3) and 4-DoF (Hopper, 0/1) never trigger onset, while 12-DoF (Quadruped) and 21-DoF (Humanoid) baselines exhibit intrinsic value overestimation (NVE $4.7\times$ and $12.97\times$ without the gate mechanism) that precludes CIVO isolation. Onset occurs exclusively at 6 DoF (Walker 60\%, Cheetah 75\%; Tables~\ref{tab:cross_env_civo} and~\ref{tab:dof_boundary}), suggesting CIVO requires both sufficient dynamic complexity for temporal embeddings to develop exploitable value bias and a stable baseline where gate-induced perturbation dominates. Practitioners deploying temporal conditioning should monitor NVE particularly in intermediate-complexity systems where baseline value estimation is stable but dynamics are complex enough for embedding co-adaptation. For transformer-based world models \citep{micheli2023iris}, whether this boundary shifts remains open. The environment-dependence extends to remedy efficacy: intervention configurations calibrated on Walker fail on Cheetah despite identical DoF, where CIVO severity is $2$--$20\times$ higher (\S\ref{sec:remedies}), underscoring that morphology-specific dynamics determine both onset risk and remedy effectiveness.

\paragraph{Independent configuration effects.}
Not all degradation in variable-rate configurations is CIVO-related; the short imagination horizon ($H{=}8$) and temporal subsampling independently reduce performance (exp\_008f: 33\% baseline without CIVO). The NVE diagnostic specifically isolates the value-overestimation component, distinguishing CIVO from these orthogonal sources of degradation.
